\cleardoublepage
\chapter*{Resumen}
\markboth{Resumen}{}
\addcontentsline{toc}{chapter}{Resumen}

El cáncer de mama sigue siendo una de las enfermedades más mortales
en la población femenina en todo el mundo, lo que hace que el
diagnóstico preciso y la planificación del tratamiento sean cruciales
para mejorar las condiciones de los pacientes. Esta tesis presenta un
enfoque novedoso para la Segmentación Semántica de Imágenes (ISS) en
imágenes médicas, con un enfoque particular en imágenes
histopatológicas de tejido de cáncer de mama. El trabajo aborda el
desafío crítico de manejar las etiquetas de múltiples anotadores en
escenarios de crowdsourcing, donde el rendimiento de los anotadores
varía a través del espacio de entrada y la calidad de la imagen. La
solución propuesta introduce una función de pérdida novedosa que
infiere implícitamente la segmentación óptima y evalúa la
confiabilidad de los etiquetadores sin requerir entradas explícitas
sobre su rendimiento. Esto se logra a través de un mapa de
confiabilidad intermedio que permite al modelo priorizar la
información de etiquetadores confiables mientras descarta etiquetas ruidosas.

La metodología se aparta de los modelos de segmentación basados en
CNN convencionales al eliminar la necesidad de supervisión directa
del rendimiento de los etiquetadores, mejorando así su generalización
y adaptabilidad a través de diversos conjuntos de datos y cohortes de
etiquetadores. El enfoque combina las fortalezas de las arquitecturas
de aprendizaje profundo en forma de U con una sofisticada función de
pérdida que modela la confiabilidad de los anotadores a través del
espacio de entrada. El rendimiento del modelo se evalúa tanto en
conjuntos de datos sintéticos como en imágenes histopatológicas del
mundo real, demostrando resultados superiores en comparación con los
enfoques más avanzados en términos de precisión de segmentación y
cuantificación de incertidumbre.

La investigación hace contribuciones significativas al campo al
proporcionar una solución robusta para manejar anotaciones
inconsistentes en la segmentación de imágenes médicas,
particularmente valiosa en histopatología donde las anotaciones de
expertos son costosas y requieren mucho tiempo para obtenerlas. El
enfoque propuesto muestra resultados prometedores en la reducción de
costos de anotación mientras mantiene una alta precisión de
segmentación, lo que lo hace particularmente relevante para
aplicaciones clínicas donde la segmentación precisa de estructuras
tisulares es crucial para el diagnóstico y la planificación del tratamiento.

\textbf{\small Palabras clave:} Crowdsourcing, Múltiples Anotadores,
Histopatología, Cáncer de Mama, Segmentación Semántica de Imágenes,
Procesos Gaussianos, Aprendizaje Profundo, Análisis de Imágenes Médicas
